\section{Strategische Ma{\ss}nahmen}
\label{sec:strategie-azn}

Tatsachenbericht. Keine Bewertung, kein Ausblick in eigener Formulierung, kein
wertendes Adjektiv. Was die Ma{\ss}nahmen taugen, steht in
Teil~\ref{sec:urteil-azn}.

\subsection{Die vier Therapiegebiete}

\begin{table}[htbp]
\centering
\caption{Umsatz nach Therapiegebiet, Mio.\,USD}
\label{tab:azn-segmente}
\begin{tabular}{lrrr}
\toprule
Therapiegebiet & 2025 & 2024 & Ver\"anderung\\
\midrule
Onkologie                        & \Betrag{\AznOnkologie} & \Betrag{\AznOnkologieVJ} & \Pct{\AznOnkologieWachstum}\\
Herz, Niere, Stoffwechsel        & \Betrag{\AznCvrm} & \Betrag{\AznCvrmVJ} & \Pct{\AznCvrmWachstum}\\
Atemwege und Immunologie         & \Betrag{\AznRespImmun} & \Betrag{\AznRespImmunVJ} & \Pct{\AznRespImmunWachstum}\\
Impfstoffe und Immuntherapien    & \Betrag{\AznImpfstoffe} & \Betrag{\AznImpfstoffeVJ} & \Pct{\AznImpfstoffeWachstum}\\
Seltene Erkrankungen             & \Betrag{\AznSeltene} & \Betrag{\AznSelteneVJ} & \Pct{\AznSelteneWachstum}\\
\bottomrule
\end{tabular}
\end{table}

Die Zahlen der Tabelle~\ref{tab:azn-segmente} stammen aus der
Therapiegebiets\"ubersicht des Berichts.\quelleGB{2025}{13}\quelleMerken{s12gbxcacfxbd}\quelleGB{2025}{17}\quelleMerken{s12gbxcacfxbh}\quelleGB{2025}{23}\quelleMerken{s12gbxcacfxcd}
Die Onkologie tr\"agt mit \MioUSD{\AznOnkologie} den gr\"o{\ss}ten und mit
\Pct{\AznOnkologieWachstum} zugleich den am schnellsten wachsenden Anteil.
Das gr\"o{\ss}te Einzelmedikament ist Farxiga mit \MioUSD{\AznFarxiga} und
damit \Pct{\AznFarxigaAnteil} des Konzernumsatzes, gefolgt von Tagrisso mit
\MioUSD{\AznTagrisso} und Imfinzi mit
\MioUSD{\AznImfinzi}.\quelleGB{2025}{18}\quelleMerken{s12gbxcacfxbi}

\subsection{Patentabl\"aufe}
\label{sec:loe-azn}

Der Bericht beziffert keinen Gesamtumsatz, der an auslaufenden Patenten
h\"angt; er verweist daf\"ur auf eine gesonderte Ver\"offentlichung.\quelleGB{2025}{30}\quelleMerken{s12gbxcacfxda}
Was er beziffert, sind einzelne F\"alle.

Brilinta verlor 2025 durch Generikaeintritt in den Vereinigten Staaten und in
Europa; der Umsatz fiel auf \MioUSD{\AznBrilinta}.\quelleErneut{s12gbxcacfxbi}
Bei Farxiga ist der Verlust der US-Marktexklusivit\"at f\"ur 2026 angek\"undigt,
zeitgleich mit der Preisverhandlung nach dem Inflation Reduction
Act.\quelleGB{2025}{33}\quelleMerken{s12gbxcacfxdd}
Im ersten Halbjahr 2026 ist er eingetreten: Farxiga kam auf
\MioUSD{\AznFarxigaHJ}, in den Vereinigten Staaten auf
\MioUSD{\AznFarxigaUsaHJ}; der Bericht nennt als Ursache
\enquote{multiple generics launched in Q2 2026} sowie die
mengenbasierte Beschaffung in China.\quelleHJ{9}\quelleMerken{s12hjxj}
Rechnet man das Halbjahr auf das Jahr hoch, liegt Farxiga
\Pct{\AznFarxigaHJWachstum} unter dem Vorjahresniveau.

\subsection{Investitionsprogramm und Zuk\"aufe}

Das Unternehmen hat ein Investitionsprogramm in den Vereinigten Staaten von
\MrdUSD{\AznInvestitionUsa} bis 2030 angek\"undigt; darin enthalten sind eine
neue Fertigungsanlage in Virginia \"uber \MrdUSD{\AznVirginia} und ein Ausbau
in Maryland \"uber \MrdUSD{\AznMaryland}.\quelleGB{2025}{4}\quelleMerken{s13gbxcacfxe}\quelleGB{2025}{33}\quelleMerken{s13gbxcacfxdd}
Im Oktober 2025 schloss das Unternehmen eine Vereinbarung mit der
US-Regierung \"uber Arzneimittelpreise, die eine dreij\"ahrige Ausnahme von
Z\"ollen nach \ZollParagraf{} einschlie{\ss}t.\quelleErneut{s13gbxcacfxdd}

An Zuk\"aufen und Lizenzen fielen 2025 an: EsoBiotec mit einer Anfangszahlung
von \MioUSD{\AznEsoBiotec}, FibroGen China mit \MioUSD{\AznFibroGen} und eine
Forschungskooperation mit CSPC \"uber eine Anfangszahlung von
\MioUSD{\AznCspc}.\quelleGB{2025}{62}\quelleMerken{s13gbxcacfxgc}\quelleGB{2025}{63}\quelleMerken{s13gbxcacfxgd}
Im Juli 2026 kam eine Lizenz von Dizal \"uber \MioUSD{\AznDizal}
hinzu.\quelleHJ{5}\quelleMerken{s13hjxf}

Die Belegschaft wuchs von \Stueck{\AznMitarbeiterVJ} auf
\Stueck{\AznMitarbeiter}, davon \Stueck{\AznMitarbeiterFuE} in Forschung und
Entwicklung.\quelleGB{2025}{27}\quelleMerken{s13gbxcacfxch}

\subsection{Umsatzambition 2030}

Das Unternehmen hat sich selbst ein Umsatzziel gesetzt, und es steht im
Wortlaut im Bericht: \enquote{By 2030, we aim to launch at least 20 new
medicines and achieve \$80 billion in Total Revenue with sustained growth
thereafter.}\quelleGB{2025}{10}\quelleMerken{s13gbxcacfxba}
Von \MioUSD{\AznUmsatzXXV} im Jahr 2025 auf \MrdUSD{\AznAmbition} bis 2030
erfordert \Pct{\AznAmbitionRate} Wachstum pro Jahr. Der Halbjahresbericht
best\"atigt die Ambition unver\"andert.\quelleHJ{1}\quelleMerken{s13hjxb}

\subsection{Prognose gegen Ist}
\label{sec:prognose-azn}

Ein Jahr Prognosetreue sagt, ob das Unternehmen seine Zahl erreicht hat. Vier
Jahre sagen, was seine Prognose wert ist, und das ist die n\"utzlichere
Auskunft. Erfasst ist jeweils die \emph{zuerst} gegebene Prognose aus dem
Bericht des Vorjahres, nicht eine sp\"ater angehobene.

\begin{table}[htbp]
\centering
\caption{Prognose, wie zuerst gegeben, gegen das Ergebnis. Alle Angaben auf
bereinigter Basis, weil das Unternehmen ausdr\"ucklich keine Prognose auf
ausgewiesener Basis abgibt.}
\label{tab:azn-prognose}
\begin{tabular}{llll}
\toprule
Jahr & Prognose Umsatz & Ist Umsatz & Lage\\
\midrule
2022 & high-teens \%              & \Pct{\AznIstXXII} & im Bereich\\
2023 & low-to-mid single-digit \% & \Pct{\AznIstXXIII} & im Bereich\\
2024 & low double-digit to low teens \% & \Pct{\AznIstXXIV} & \"uber dem Bereich\\
2025 & high single-digit \%       & \Pct{\AznIstXXV} & im Bereich\\
\bottomrule
\end{tabular}
\end{table}

Die Prognosen stehen im jeweiligen Vorjahresbericht,\quelleGB{2021}{66}\quelleMerken{s14gbxcacbxgg}\quelleGB{2022}{72}\quelleMerken{s14gbxcaccxhc}\quelleGB{2023}{71}\quelleMerken{s14gbxcacdxhb}\quelleGB{2024}{81}\quelleMerken{s14gbxcacexib}
die Ergebnisse im Bericht des Jahres selbst.\quelleGB{2022}{60}\quelleMerken{s14gbxcaccxga}\quelleGB{2023}{58}\quelleMerken{s14gbxcacdxfi}\quelleGB{2024}{67}\quelleMerken{s14gbxcacexgh}\quelleGB{2025}{55}\quelleMerken{s14gbxcacfxff}
F\"ur 2021 ist die zuerst gegebene Prognose nicht belegbar: Sie st\"unde im
Bericht 2020, dessen Seitenzahlen nicht auslesbar sind. Die Reihe umfasst
deshalb vier statt f\"unf Jahre, und diese L\"ucke wird nicht durch eine
sp\"ater ver\"offentlichte Fassung gef\"ullt.

In keinem der vier Jahre lag das Ergebnis unter der zuerst gegebenen Spanne;
in einem lag es dar\"uber, und im Bericht steht dazu, dass die Prognose
w\"ahrend des Jahres zweimal angehoben wurde.\quelleErneut{s14gbxcacexgh}

F\"ur 2026 lautet die Prognose, unver\"andert seit Februar 2026 und im Juli
best\"atigt: Umsatzwachstum im mittleren bis hohen einstelligen
Prozentbereich, Wachstum des bereinigten Ergebnisses je Aktie im niedrigen
zweistelligen Bereich.\quelleHJ{1}\quelleMerken{s14hjxb}
Das erste Halbjahr lieferte \MioUSD{\AznUmsatzHJ} Umsatz und
\USDje{\AznCoreErgebnisJeAktieHJ} bereinigtes Ergebnis je
Aktie.\quelleErneut{s14hjxb}
